
\section{挑战与未来展望}
\label{sec:challenges}

\subsection{当前面临的主要挑战}
\label{sec:open_issues}

\noindent\textbf{人群差异：}血管弹性、外周阻力、皮肤色素与佩戴位置因人而异，同一 PPG 波形在不同人群中的血压映射并不一致；在手术、ICU 等特定人群上训练的模型，迁移到健康人群或跨设备场景时误差明显上升。

\noindent\textbf{个性化校准：}腕部 PPG 与上臂袖带血压的关系受外周血管与姿势影响，通常需要基线校准；单点校准难以覆盖昼夜与状态变化，频繁校准又影响可穿戴体验，如何利用个人历史数据中的稳态信息、以尽量少的校准跟踪个体漂移，尚无统一方案。

\noindent\textbf{标签稀缺：}可靠的血压标签依赖袖带或动脉线，采集成本高，公开数据或来自手术与 ICU、或为静息短时记录，腕部带标签数据尤其稀少；标签还含测量误差与时间不同步，而无标签 PPG 虽易获取，却难以直接用于监督学习。

\subsection{未来研究方向}
\label{sec:future}

针对上述挑战，可以从三个方向展开研究。其一，跨人群泛化：按年龄、肤色、体位等分层评测，结合域自适应与个性化微调，降低对特定人群的依赖。其二，低负担校准：利用个人历史片段中的稳态信息与少样本学习，在单点甚至零校准条件下完成适配。其三，标签高效学习：以自监督预训练利用海量无标签 PPG，结合弱标签与迁移学习。
